% Generated draft: requires mathematical and source review.
\documentclass[11pt]{article}
\usepackage[margin=25mm]{geometry}
\usepackage{fontspec}
\setmainfont{Latin Modern Roman}
\usepackage{amsmath,amssymb,mathtools,unicode-math}
\setmathfont{Latin Modern Math}
\usepackage{enumitem,needspace,xurl,hyperref}
\setlength{\parindent}{0pt}
\setlength{\parskip}{5pt}
\setlength{\emergencystretch}{3em}
\newcommand{\N}{\mathbb{N}}
\newcommand{\Z}{\mathbb{Z}}
\newcommand{\Q}{\mathbb{Q}}
\newcommand{\R}{\mathbb{R}}
\newcommand{\C}{\mathbb{C}}
\newcommand{\F}{\mathbb{F}}
\newcommand{\A}{\mathbb{A}}
\newcommand{\PP}{\mathbb{P}}
\newcommand{\E}{\mathbb{E}}
\newcommand{\eps}{\varepsilon}
\newcommand{\dd}{\,\mathrm{d}}
\begin{document}
\section*{UnsolvedMath ID: 30005316}
\textbf{Source problem code:} catalogue entry
\subsection{Definitions and mathematical statement}
Let a partition be a finite weakly decreasing sequence \(\lambda=(\lambda_1,\ldots,\lambda_\ell)\) of positive integers, with the empty partition also allowed. Its Young diagram is \(D_\lambda=\{(i,j):1\le i\le\ell,\ 1\le j\le\lambda_i\}\). A domino consists of two cells sharing an edge. A partition has empty \(2\)-core precisely when its Young diagram can be tiled by dominoes. The supplied excerpt gives \((6,4,3,1)\) as an example. Write \(P_\lambda(x;q,t)\) for the Macdonald symmetric polynomials, indexed by partitions, with parameters \(q,t\); their Schur specialization is \(q=t\).

\textbf{Requested conjectural assertion, specified by reference.} Establish the pair of vanishing-integral identities for Macdonald polynomials attributed to Lee, Rains and Warnaar in reference [3, Conjecture 9.2] of the supplied Oberwolfach excerpt, throughout their intended parameter ranges, including the general case beyond the Schur specialization \(q=t\).

\textbf{Unresolved formulation.} The supplied record does not reproduce either integral or its proposed evaluation. It also omits the integration variables, measure, domain or contour, auxiliary parameters, admissibility and convergence hypotheses, partition restrictions, and exact vanishing conditions. Thus the assertion above identifies the intended target by reference but is not a self-contained, fully quantified mathematical conjecture. The excerpt associates the surrounding conjectures with partitions having empty \(2\)-core; it does not specify how that condition enters either of the two integral identities. No vanishing criterion can therefore be inferred from that association alone. The exact identities and their hypotheses must be supplied before a precise standalone statement can be completed.
\subsection{Short English statement}
The question asks whether the particular pair of Lee--Rains--Warnaar vanishing-integral conjectures holds for general Macdonald parameters beyond \(q=t\). The supplied curation credits Albion with proving the Schur case. The actual integral formulas and hypotheses are missing, so this draft identifies the target but requires source recovery and human mathematical review.
\subsection{Sources}
\begin{itemize}
\item Local UnsolvedMath registry, supplied source record: permanent ID \(30005316\), exact source problem code \texttt{catalogue entry}, title ``Lee--Rains--Warnaar Macdonald Integral Identities.'' The record gives proposed year 2022 and no named proposer.
\item Supplied report citation: \emph{Enumerative Combinatorics} (2023), Oberwolfach Reports, pp.\ 3215--3305, DOI \texttt{10.4171/owr/2022/57}. \url{https://doi.org/10.4171/owr/2022/57}. This is also the supplied literature source labelled ``original OWR report.''
\item Supplied literature source labelled ``Albion primary article'': \url{https://doi.org/10.1090/bproc/221}. The supplied assessment credits Albion with proving the Schur specialization and describes the general Macdonald case of the extracted conjecture as unresolved.
\item Literature assessment date: \texttt{2026-08-22}. This is supplied dated curation, not an independently verified literature assessment. The supplied status is \texttt{partially\_solved}.
\item The excerpt attributes the target pair to Lee, Rains and Warnaar, reference [3, Conjecture 9.2], and attributes a proof of the Schur case to its reference [1]. It also mentions the virtual Koornwinder integral approach to Littlewood identities of Rains and Warnaar, reference [6], and a connection with MacMahon's conjecture for the generating function of symmetric plane partitions in an \(n\times n\times m\) box. Full bibliographic entries for [1], [3], and [6] are absent.
\item The registry labels the statement \texttt{corrected\_verified} and reports that the official report was checked and the question isolated with repaired notation. These are supplied provenance claims; no independent source inspection is asserted here.
\end{itemize}
\end{document}
